% ===================================================================
%  تابلوی شمارهٔ ۷ — روستا در زمستان و خانهٔ روستایی در بهار.
%
%  Traduction, faite en 2026, en persan d'aujourd'hui, sur l'ido de
%  texto/io. Regles de la colonne : voir 10-tablo-01.tex. Deux
%  scenes, deux numerotations. Ouverture de la DEUXIEME SERIE :
%  « سری دوم ». Un bloc marque « ¶2 » par ordo.py, t07-c1-07-1 : on
%  n'en ouvre pas une seconde cle.
%
%  LA FERME EST LE QUATRIEME LEXIQUE ENTIEREMENT IRANIEN DE CETTE
%  COLONNE, APRES LE CORPS, LES SENS ET LA PARENTE — ET C'EST LE PLUS
%  ANCIEN DE TOUS. گاو la vache (25), اسب le cheval (7), گوسفند le
%  mouton (50), بز la chevre (51), مرغ la poule (43), خروس le coq
%  (71), غاز l'oie (69), اردک le canard (66), سگ le chien (37), خر
%  l'ane (31), گندم le ble, کاه la paille, یونجه la luzerne, خرمن
%  l'aire, داس la faucille, خیش la charrue (5), شیار le sillon (8),
%  چاه le puits (29), آخور la mangeoire, طویله l'etable (58), کندو la
%  ruche (21), شیر le lait, کره le beurre, پنیر le fromage (64).
%
%  ET PLUSIEURS DE CES MOTS SONT INDO-EUROPEENS SANS INTERMEDIAIRE :
%  گاو est le meme mot que le latin bos, l'anglais cow et le sanskrit
%  go ; اسب est le latin equus ; مادر, پدر, برادر du tableau 4
%  etaient de la meme famille. Le persan est la seule colonne du
%  releve qui partage une racine avec le francais SANS l'avoir
%  empruntee — et le tableau de la ferme est celui ou cela se voit le
%  mieux, parce que ce sont les mots qu'on n'emprunte jamais.
%
%  LA NEIGE (37) EST ICI CHEZ ELLE, ET LA COMPARAISON AVEC JAVA EST
%  ENTIERE. Le tableau 7 javanais avait du emprunter « salju » pour la
%  neige, « es » pour la glace, et composer « sepatu es » pour le
%  patin (26) — trois procedes pour une saison qui n'existe pas la-
%  bas. Le persan ecrit برف, یخ et اسکیت : deux mots iraniens
%  ordinaires et un seul emprunt, pour le seul objet qui soit
%  reellement venu d'ailleurs. La langue emprunte le PATIN, jamais la
%  GLACE.
%
%  ET LE MOULIN (16) TRANCHE AUTREMENT QU'A JAVA. Java n'a jamais eu
%  de moulin a vent et le javanais a du composer « kincir angin » ;
%  l'Iran a invente le moulin a vent — les آسبادهای نشتیفان du
%  Khorassan tournent depuis le IXe siecle et sont les plus anciens du
%  monde. آسیاب est donc un mot ordinaire, et آسیاب بادی n'est pas une
%  traduction mais un nom. La meme planche, le meme objet, et deux
%  langues dont l'une l'a recu et l'autre l'a fait.
%
%  UNE FAUTE D'ORDRE, ET C'EST ENCORE UN COMPLEMENT MIS DU MAUVAIS
%  COTE — la troisieme du meme genre dans cette colonne, apres les
%  deux du tableau 6. Au bloc t07-c2-15-1 le fumier (81) avait ete
%  pose apres les poussins (75), alors que l'ido dit « marodas sur la
%  sterko kun sua yuneti ». Le persan met ses complements dans
%  l'ordre qu'il veut, et c'est justement pourquoi la main s'y trompe.
%  Trois fautes, trois fois le meme mecanisme, et zero faute
%  structurelle : cette colonne n'a aucune difficulte de grammaire,
%  elle n'a qu'une difficulte d'attention.
%
%  LE VILLAGE EMPRUNTE, LUI, ET AU FRANCAIS COMME PARTOUT DANS CETTE
%  COLONNE : دلیجان la diligence (11), بیمه l'assurance — non, celui-ci
%  est persan —, پست, بلیت. Mais دلیجان est le mot le plus curieux du
%  tableau : il vient du francais « diligence », et il est entre en
%  persan au XIXe siecle avec la voiture elle-meme, exactement comme
%  les mois gregoriens du tableau 3. Un vehicule et un calendrier
%  arrives par la meme porte.
% ===================================================================

%%K t07-noto-1 noto
(*) در بارهٔ جزئیات خوراک، تابلوهای ۶ و ۱۳ را ببینید.

%%K t07-apar-1 apar
\VUsaut{4.0mm}
\VUcentre{13.2pt}{90}{سری دوم}
\VUsaut{3.0mm}
\VUfilet{16mm}
\VUsaut{5.6mm}
\VUtitre{15.0pt}{60}{تابلوی شمارهٔ ۷}
\VUsaut{3.0mm}

%%K t07-tit-1 sub
\VUsaut{3.0mm}
\VUcentre{12.0pt}{0}{\textit{صحنهٔ اول.}}
\VUsaut{3.0mm}

%%K t07-tit-2 sub
\VUsaut{4.0mm}
\VUpk{10.2pt}{40}{روستا در زمستان.}
\VUsaut{4.0mm}

%%K t07-c1-01-1 p
1. --- در تعطیلات سال نو، پدرم را تا نووُربا همراهی کردم، روستای کوچکی
که در آنجا باید چند کار بازرگانی را سامان می‌داد.

%%K t07-c1-02-1 p
2. --- ما با \VUgras{دلیجان} \textsuperscript{(11)} رفتیم که میان شهر و
آن قصبه کار می‌کند ؛ اما چون بسیار سرد بود، این سفر با گردونه‌ای کم‌آسایش
چندان خوشایند نبود.

%%K t07-c1-03-1 p
3. --- از این رو شادی راستین یافتیم هنگامی که دیدیم \VUgras{کالسکه‌ران}
گردونه‌اش را در میدان، جلوِ \VUgras{مهمانخانهٔ} \textsuperscript{(2)}
\textit{خرس سپید} نگه داشت. \VUgras{مهمانخانه‌دار}
\textsuperscript{(4)} بر آستانهٔ خانه‌اش در انتظار ما بود و آرام
\VUgras{پیپ} خود را دود می‌کرد.

%%K t07-c1-03-2 p
او ما را به درون خواند و من درنگ نکردم و جلوِ آتش شاخه‌های خشک نشستم که
در اجاق تق‌تق می‌کرد. آنجا \VUgras{باد شمال} تند را دست انداختم که
\VUgras{تابلوی} \textsuperscript{(3)} کهنه را به صدا درمی‌آورد و
\VUgras{دود} \textsuperscript{(1)} بیرون آمده از دودکش را در گرداب‌های
سیاه با خود می‌برد.

%%K t07-c1-04-1 p
4. --- هنگام ناهار بود. بی‌آنکه بیشتر درنگ کنیم، خوراک ساده اما
خوشمزه‌ای را که برایمان آوردند خوب خوردیم (*).

%%K t07-tit-3 sub
\VUsaut{3.0mm}
\VUpk{10.2pt}{40}{چند دیدار.}
\VUsaut{3.0mm}

%%K t07-c1-05-1 p
5. --- پس از ناهار، پدرم را که نمایندهٔ \VUgras{بیمه} است تا نزد
مشتریانش همراهی کردم. نخست نزد \VUgras{آهنگر} \textsuperscript{(5)}
رفتیم که نزدیک مهمانخانه می‌نشیند. او سرگرم \VUgras{نعل} زدن به اسبی
بود. از نیروی همکارش شگفت‌زده شدم که \VUgras{سم} اسب را استوار بر
پیش‌بند چرمی خود نگه داشته بود، در حالی که آهنگر \VUgras{نعل}
\textsuperscript{(6)} را با ضربه‌های نیرومند چکش می‌کوبید.

%%K t07-c1-06-1 p
6. --- سپس نزد \VUgras{کفشگر} \textsuperscript{(7)} روستا رفتیم،
یاکوبوس پیر. با آنکه \VUgras{چکمهٔ} بسیار زیبایی را تابلوی خود کرده
است، به کفی‌زدن یک جفت کفش کهنهٔ سوراخ خرسند است.

%%K t07-c1-06-2 p
پیرمرد خندان به ما گفت : « در شهر، من « کفش‌دوز » بودم و مشتری
نداشتم. اینجا، من « کفش‌دوزِ پینه‌دوز » هستم و کار بسیار است. »

%%K t07-c1-07-1 p
7. --- او از همسایه‌اش، \VUgras{سلمانی} \textsuperscript{(8)}، برایمان
گفت که مشتریانش خرسند نیستند. \VUgras{تیغ‌هایش} همیشه لب‌پریده‌اند و
\VUgras{قیچی‌هایش} هرگز خوب نمی‌برند.

اما چون او تنها سلمانی روستاست، مردم ناچارند نزد او ریش بتراشند و
موی خود را کوتاه کنند. وانگهی او مردی خسیس و ناخوشایند است.

%%K t07-c1-08-1 p
8. --- خوشبختانه، \VUgras{همسایگان} دیگر مانند او نیستند. برای نمونه
\VUgras{گاری‌ساز} \textsuperscript{(9)} مردی مهربان است، کوشا و
یاری‌رسان. تعمیر گردونه نزد او خوشایند است. کارش همیشه به پایان رسیده
است.

%%K t07-c1-09-1 p
9. --- یاکوبوس پیر از \VUgras{زین‌ساز} \textsuperscript{(10)} هم گله
نداشت، که گاهی هنگام فشار کار در دوختن \VUgras{یراق} به او یاری
می‌دهد.

%%K t07-c1-09-2 p
پدرم از خانواده‌اش پرسید : « همه خوب‌اند. نوهٔ دخترم در
\VUgras{مدرسه} \textsuperscript{(12)} است. \VUgras{آموزگارش}
\textsuperscript{(13)} از او بسیار خرسند است، او \VUgras{دانش‌آموز}
\textsuperscript{(14)} خوبی است. مانند پسر شیطان من نیست ؛ او تنها به
بازی می‌اندیشد. \VUgras{آموزگار} \textsuperscript{(15)} نمی‌تواند چیزی
به او بیاموزد. »

%%K t07-tit-4 sub
\VUpk{10.2pt}{40}{گپ‌های روستایی.}
\VUsaut{3.0mm}

%%K t07-c1-10-1 p
10. --- پیرمرد سپس خبرهای روستا را برایمان گفت. مدرسه‌ای نو خواهند
ساخت، زیرا آنکه در \VUgras{خانهٔ ده} برپاست بسیار کوچک شده است. وانگهی
این آسان است، زیرا خریدن بخشی از باغ \VUgras{آسیابان} بس خواهد بود.
این برای آن مرد بینوا بسیار سودمند خواهد بود. از سرما،
\VUgras{سنگ‌های آسیاب} \VUgras{آسیاب} \textsuperscript{(16)} دیگر
نمی‌توانند گندم را آرد کنند، زیرا \VUgras{چرخ} \textsuperscript{(17)}
را \VUgras{یخ} \textsuperscript{(25)} \VUgras{جویبار}
\textsuperscript{(23)} بسته است.

%%K t07-c1-11-1 p
11. --- « در بارهٔ ساختمان، آیا \VUgras{کلیسای} \textsuperscript{(18)}
نو ما را دیدید؟ زیر ردای برفی خود بسیار زیباست. \VUgras{کلاغ‌ها}
\textsuperscript{(19)} که دیگر ناقوس‌خانهٔ کهنهٔ خود را نمی‌شناسند،
غمگین بر درخت بزرگ \VUgras{گورستان} \textsuperscript{(20)}
می‌نشینند. »

%%K t07-c1-12-1 p
12. --- این اندیشه در بارهٔ گورستان کفشگر را به یاد کسانی انداخت که
پدرم می‌شناخت و سال گذشته درگذشتند. « چه بسیار \VUgras{گور}
\textsuperscript{(21)}، آقای عزیز، زیر \VUgras{سرو}
\textsuperscript{(22)} به دیگران افزوده شد! مرگ دفتردار پیر ما را درو
کرد. همهٔ روستاییان در \VUgras{خاکسپاری} او بودند. »

%%K t07-c1-13-1 p
13. --- « این هم جانشین او که روی جویبار اسکیت می‌کند. همین حالا آمد تا
تسمهٔ \VUgras{اسکیت} \textsuperscript{(26)} خود را که پاره شده بود نزد
من درست کند. »

%%K t07-c1-14-1 p
14. --- « این هم بر کرانهٔ جویبار، \VUgras{نگهبان مزارع}
\textsuperscript{(27)} نو. او ژاندارم پیشین است که بچه‌های شیطان روستا
را می‌ترساند. آنان همین که \VUgras{ساق‌پوش} \textsuperscript{(29)} و
\VUgras{کلاه کپی} \textsuperscript{(28)} او را ببینند، می‌گریزند. »

%%K t07-c1-15-1 p
15. --- « آه! این هم یوسف پیر که \VUgras{گاری کوچک هیزم}
\textsuperscript{(30)} را برای نانوا می‌برد. --- درست است، پدرم گفت،
\VUgras{قاطرهای} \textsuperscript{(31)} او را می‌شناسم. به راستی باید با
او سخن بگویم ؛ از این فرصت بهره خواهم برد. تا دیدار، پدر یاکوبوس. »

%%K t07-c1-16-1 p
16. --- درست هنگامی که بیرون آمدیم، کودکان روستا مدرسه را ترک کردند و
دوان‌دوان به \VUgras{سرسرهٔ یخ} \textsuperscript{(33)} شتافتند، با خطر
افتادن و شکستن اندامی در \VUgras{زمین‌خوردن} \textsuperscript{(34)}.
یکی از آنان \VUgras{سورتمهٔ} \textsuperscript{(32)} خود را از گاری‌ساز
گرفت، یکی از رفیقانش را در آن نشاند و دوان به راه افتاد.

%%K t07-c1-17-1 p
17. --- دیگران به \VUgras{تودهٔ برف} \textsuperscript{(37)} نزدیک
\VUgras{پل} \textsuperscript{(40)} شتافتند و \VUgras{آدم برفی}
\textsuperscript{(39)} را بمباران کردند که دیروز برپا کرده بودند و کلاه،
پیپ و کیف مدرسه بر او گذاشته بودند.

%%K t07-c1-18-1 p
18. --- آنان به باد یخبندان و سرما چندان اهمیت نمی‌دهند. بیشترشان حتی
\VUgras{شال گردن} \textsuperscript{(35)} ندارند، و برای گرم شدن پس از
نبرد \VUgras{گلوله‌های برفی} \textsuperscript{(38)}، مشتی \VUgras{شاه‌بلوط}
\textsuperscript{(42)} بسیار داغ بس است که از \VUgras{فروشندهٔ} پیر،
یاکوبینا، خریده‌اند که زیر \VUgras{شال} \textsuperscript{(43)} بسیار
نازک خود از سرما می‌لرزد.

%%K t07-tit-5 sub
\VUsaut{3.0mm}
\VUcentre{12.0pt}{0}{\textit{صحنهٔ دوم.}}
\VUsaut{3.0mm}

%%K t07-tit-6 sub
\VUsaut{4.0mm}
\VUpk{10.2pt}{40}{خانهٔ روستایی در بهار.}
\VUsaut{4.0mm}

%%K t07-c2-01-1 p
1. --- با همهٔ خوشی‌های زمستان، بازگشت \VUgras{بهار} را با شادی ژرف
خوشامد می‌گوییم.

%%K t07-c2-01-2 p
حیاط خانهٔ روستایی، شامگاهان، در ماه مه، چه تابلوی شادی‌بخشی می‌نماید!

%%K t07-c2-02-1 p
2. --- در افق \VUgras{خورشید} \textsuperscript{(1)} آرام فرو می‌نشیند و
\VUgras{دشت} \textsuperscript{(3)} را با \VUgras{پرتوهای}
\textsuperscript{(2)} ارغوانی خود غرق می‌کند. \VUgras{شخم‌زن}
\textsuperscript{(6)} کوشا می‌شتابد تا \VUgras{کشتزار}
\textsuperscript{(4)} خود را شخم بزند.

%%K t07-c2-02-2 p
او با \VUgras{خیش} \textsuperscript{(5)} خود که به \VUgras{اسب}
\textsuperscript{(7)} نیرومندی بسته شده است، \VUgras{شیار}
\textsuperscript{(8)} می‌کشد که \VUgras{بذرافشان} \textsuperscript{(9)}
در آن با مشت پر \VUgras{بذر} خواهد پاشید که خوشه‌های زرین خواهد داد.

%%K t07-c2-03-1 p
3. --- \VUgras{دروگران} \textsuperscript{(11)} با داس‌های خود علف مرتع
را درویده‌اند، خشک‌کنندگان \VUgras{یونجه} \textsuperscript{(10)} را با
گرداندنش با \VUgras{چنگک} \textsuperscript{(12)} و شن‌کش خشک کرده‌اند،
و این هم که بازمی‌گردند.

%%K t07-c2-03-2 p
برخی جلوِ گردونهٔ سنگین که گاوها آن را می‌کشند راه می‌روند ؛ ابزار خود
را بر دوش می‌برند. دیگران، تنبل‌تر، آسوده روی یونجهٔ خوشبو نشسته‌اند.

%%K t07-c2-04-1 p
4. --- هنگام گذشتن، به \VUgras{باغبان} \textsuperscript{(16)} سلام
دوستانه می‌دهند که در \VUgras{باغ میوه} \textsuperscript{(13)} سرگرم
هرس \VUgras{درختان میوه} و پیوند زدن \VUgras{درختان گلابی}
\textsuperscript{(14)} است. \VUgras{درختان آلو} \textsuperscript{(15)}
شکوفه می‌دهند، و در \VUgras{باغ سبزی} \textsuperscript{(17)} که خوب
کود داده شده است، سبزی‌ها، کلم‌ها و کاهوها از پیش خوب‌اند.

%%K t07-c2-04-2 p
روی دیوار کوتاه، \VUgras{طاووس} \textsuperscript{(18)} زیبایی دم خود را
باز می‌کند تا پرهای بسیار زیبایش را به تماشا بگذارد.

%%K t07-tit-7 sub
\VUsaut{3.0mm}
\VUpk{10.2pt}{40}{حیاط خانهٔ روستایی.}
\VUsaut{3.0mm}

%%K t07-c2-05-1 p
5. --- درهای \VUgras{اسطبل} \textsuperscript{(19)} باز است و آخور و
\VUgras{بستر} \textsuperscript{(23)} \VUgras{مادیانی}
\textsuperscript{(21)} دیده می‌شود که \VUgras{مهتر}
\textsuperscript{(20)} همین حالا آن را از یوغ باز کرده است.
\VUgras{کره‌اسب} \textsuperscript{(22)} که با پاهای بلند خود چنین
بامزه است، جست‌وخیزکنان مادرش را دنبال می‌کند.

%%K t07-c2-06-1 p
6. --- نزدیک \VUgras{چاه} \textsuperscript{(29)}، \VUgras{گاوها}
\textsuperscript{(25)} می‌روند تا در \VUgras{آبشخور چوبی}
\textsuperscript{(26)} که برای آنان آبشخور است بنوشند.
\VUgras{گوساله} \textsuperscript{(24)} از این فرصت بهره می‌برد تا شیر
بمکد، و \VUgras{گاو نر} \textsuperscript{(27)} که بوی خوش انبار علف را
می‌شنود، شادمانه می‌غرد و سر خود را با \VUgras{شاخ‌های}
\textsuperscript{(28)} نیرومند به گردنی بلند می‌کند.

%%K t07-c2-07-1 p
7. --- همه کار می‌کنند. \VUgras{خدمتکار} \textsuperscript{(32)}
\VUgras{طناب} \textsuperscript{(31)} چاه را به دست دارد. او با
\VUgras{سطل} \textsuperscript{(30)} آب می‌کشد تا به چارپایان بنوشاند.
نزدیک \VUgras{انبار غله} \textsuperscript{(33)}، مردی کاه‌ها را شن‌کش
می‌کشد، و جلوِ درِ \VUgras{خانه} \textsuperscript{(34)} \VUgras{زن
ریسندهٔ} \textsuperscript{(35)} پیری با چرخ ریسندگی و \VUgras{دوک}
خود \VUgras{کتان} یا \VUgras{شاهدانه} می‌ریسد، مانند روزگار کهن.

%%K t07-tit-8 sub
\VUsaut{3.0mm}
\VUpk{10.2pt}{40}{کشاورز.}
\VUsaut{3.0mm}

%%K t07-c2-08-1 p
8. --- \VUgras{کشاورز} \textsuperscript{(36)} همهٔ این کارها را
سرپرستی می‌کند و به خدمتکارانش دستور می‌دهد. او سفارش می‌کند که
\VUgras{ماله} \textsuperscript{(40)} و \VUgras{کلنگ}
\textsuperscript{(39)} را که نزدیک \VUgras{لانهٔ} \textsuperscript{(38)}
\VUgras{سگ} \textsuperscript{(37)} فراموش کرده‌اند زیر سقف ببرند.

%%K t07-c2-09-1 p
9. --- فریادهای او \VUgras{کبوتری} \textsuperscript{(41)} را می‌ترساند
که با همهٔ بال به \VUgras{کبوترخانه} \textsuperscript{(42)} می‌پرد، و
\VUgras{مرغ‌هایی} \textsuperscript{(43)} را که می‌شتابند تا دوباره به
\VUgras{مرغدانی} \textsuperscript{(44)} درآیند.

%%K t07-c2-10-1 p
10. --- جلوِ \VUgras{آغل گوسفند} \textsuperscript{(48)}، این هم
\VUgras{چوپان} \textsuperscript{(45)} پیر که \VUgras{گلهٔ}
\textsuperscript{(49)} خود را بازمی‌گرداند. با \VUgras{چوبدستی}
\textsuperscript{(46)} خود که هرگز از او جدا نمی‌شود، مانند
\VUgras{بلوز} \textsuperscript{(47)} رنگ‌باخته‌اش، \VUgras{قوچ‌ها}
\textsuperscript{(50)}، \VUgras{میش‌ها} \textsuperscript{(53)}،
\VUgras{بره‌ها} \textsuperscript{(54)}، \VUgras{بزها}
\textsuperscript{(51)} و \VUgras{بزغاله‌ها} \textsuperscript{(52)} را
به آغل می‌برد. \VUgras{سگ} \textsuperscript{(55)} وفادارش، آرگوس،
واپسین می‌آید و با پارس خود کندروان را به شتاب وامی‌دارد.

%%K t07-c2-11-1 p
11. --- این همه هیاهو و جنبش آرامش \VUgras{گاو شیرده}
\textsuperscript{(56)} خوب را برهم نمی‌زند که \VUgras{زنگولهٔ}
\textsuperscript{(57)} خود را به صدا درمی‌آورد، در حالی که او را جلوِ
درِ \VUgras{طویله} \textsuperscript{(58)} می‌دوشند.

%%K t07-c2-12-1 p
12. او گویی درمی‌یابد که باید شیر خود را بدهد تا \VUgras{زن کشاورز}
\textsuperscript{(60)} بتواند در \VUgras{کره‌گیر} \textsuperscript{(61)}
\VUgras{کره} آماده کند یا \VUgras{پنیرهای} \textsuperscript{(64)} بسیار
خوبی که از تختهٔ گذاشته بر \VUgras{تشت بزرگ} \textsuperscript{(63)}
می‌چکند.

%%K t07-c2-13-1 p
13. --- همهٔ \VUgras{شیرخانه} \textsuperscript{(59)} را \VUgras{کارگر
کشاورزی} \textsuperscript{(65)} بسیار پاکیزه نگه می‌دارد که همین حالا
\VUgras{ظرف‌های شیر} \textsuperscript{(62)} را در سطل بزرگی که به دست
دارد شسته است.

%%K t07-tit-9 sub
\VUsaut{3.0mm}
\VUpk{10.2pt}{40}{مرغدانی.}
\VUsaut{3.0mm}

%%K t07-c2-14-1 p
14. --- با کفش‌های چوبی کلفت خود، اندکی سنگین راه می‌رود، و چنین
\VUgras{بوقلمون} \textsuperscript{(68)}، \VUgras{اردک‌های ماده}
\textsuperscript{(66)}، \VUgras{اردک‌های نر} \textsuperscript{(67)} و
\VUgras{غازها} \textsuperscript{(69)} را می‌پراکند که \VUgras{نوک}
\textsuperscript{(70)} خود را گشاد باز می‌کنند و آشفته فریاد می‌کشند.

%%K t07-c2-14-2 p
\VUgras{خروس} \textsuperscript{(71)} که ستیزه‌جوتر است، \VUgras{تاج}
\textsuperscript{(72)} سرخ خود را راست می‌کند، پرهای \VUgras{دم}
\textsuperscript{(73)} خود را می‌تکاند و « قوقولی‌قوقو » ی پرطنینی سر
می‌دهد.

%%K t07-c2-15-1 p
15. --- \VUgras{مرغ مادر} \textsuperscript{(74)} روی \VUgras{کود}
\textsuperscript{(81)} با \VUgras{جوجه‌هایش} \textsuperscript{(75)} دانه
می‌چیند. \VUgras{بچه‌خوک} \textsuperscript{(78)} به سوی مادرش
می‌گریزد، \VUgras{خوک ماده} \textsuperscript{(77)} بسیار بزرگی که
نزدیک خوکدانی می‌بینیم.

%%K t07-c2-16-1 p
16. --- در قفس‌های خود، \VUgras{خرگوش‌ها} \textsuperscript{(79)}
\VUgras{علف} \textsuperscript{(80)} لطیفی را که دختر کشاورز برایشان
می‌آورد با ولع می‌خورند. یوهانینا خرگوش‌هایش را بسیار دوست دارد و هر
روز، پس از برداشتن \VUgras{تخم‌مرغ‌های} \textsuperscript{(76)} تازه از
مرغدانی، به دیدار دوستان کوچک خود می‌آید.

\VUsaut{6.0mm}
\VUfilet{22mm}
